%==============================================================
%  resume.tex — DevOps / Red Team Engineer Resume
%  Compiled automatically by GitHub Actions on every push
%  Edit this file → push → PDF auto-releases
%==============================================================

\documentclass[10pt, letterpaper]{article}

% ── Packages ─────────────────────────────────────────────────
\usepackage[
  top=0.5in, bottom=0.5in,
  left=0.6in, right=0.6in
]{geometry}
\usepackage{enumitem}
\usepackage{fontawesome5}
\usepackage{hyperref}
\usepackage{titlesec}
\usepackage{xcolor}
\usepackage{multicol}
\usepackage{tabularx}
\usepackage[T1]{fontenc}
\usepackage{lmodern}

% ── Colors ───────────────────────────────────────────────────
\definecolor{accent}{HTML}{89b4fa}      % Catppuccin blue
\definecolor{darktext}{HTML}{1e1e2e}
\definecolor{mutedtext}{HTML}{6c7086}
\definecolor{rulecolor}{HTML}{313244}

% ── Hyperlink style ──────────────────────────────────────────
\hypersetup{
  colorlinks=true,
  urlcolor=accent,
  linkcolor=accent
}

% ── Section formatting ───────────────────────────────────────
\titleformat{\section}
  {\large\bfseries\color{darktext}}
  {}{0em}{}
  [\color{rulecolor}\titlerule]
\titlespacing{\section}{0pt}{8pt}{4pt}

% ── No paragraph indent ──────────────────────────────────────
\setlength{\parindent}{0pt}
\setlength{\parskip}{0pt}

% ── List spacing ─────────────────────────────────────────────
\setlist[itemize]{
  leftmargin=1.2em,
  itemsep=1pt,
  parsep=0pt,
  topsep=2pt
}

%==============================================================
\begin{document}
\pagestyle{empty}

%── Header ────────────────────────────────────────────────────
\begin{center}
  {\LARGE\bfseries YOUR NAME}\\[4pt]
  {\color{mutedtext}\small
    \faIcon{map-marker-alt}~Panvel, Maharashtra, India \quad
    \faIcon{envelope}~\href{mailto:you@email.com}{you@email.com} \quad
    \faIcon{phone}~+91-XXXXXXXXXX \quad
    \faIcon{github}~\href{https://github.com/YOUR_USERNAME}{github.com/YOUR\_USERNAME} \quad
    \faIcon{linkedin}~\href{https://linkedin.com/in/YOUR_USERNAME}{linkedin.com/in/YOUR\_USERNAME}
  }
\end{center}

\vspace{2pt}

%── Summary ───────────────────────────────────────────────────
\section{Summary}
DevOps Engineer with experience in managing large-scale hybrid infrastructure (on-prem + AWS), supporting 10,000+ Linux servers across globally distributed data centers. Experienced in CI/ CD automation with Jenkins and Ansible, Infrastructure as Code using Terraform, and secure DevSecOps integrations. Proven track record of enabling reliable weekly production releases using canary deployments and high-availability architectures. CEH-certified with a strong focus on security-driven automation

%── Skills ────────────────────────────────────────────────────
\section{Technical Skills}

\begin{tabularx}{\textwidth}{@{}lX@{}}
  \textbf{Cloud \& Infra}   & AWS (EC2, IAM, S3, EKS, VPC, Auto Scaling, ALB), Terraform, Linux \\[2pt]
  \textbf{IaC \& Automation} & Bash, Python, YAML \\[2pt]
  \textbf{Containers}        & Docker, Kubernetes, Helm \\[2pt]
  \textbf{CI/CD}             & Jenkins, GitHub Actions, GitLab CI/ CD, Ansible, Git, SonarQube \\[2pt]
  \textbf{Observability}     & Prometheus, Grafana, ELK Stack, Datadog \\[2pt]
  \textbf{Security / Red Team} & Nmap, Metasploit, Burp Suite, BloodHound, Cobalt Strike, Nuclei, OSINT \\[2pt]
\end{tabularx}

%── Experience ────────────────────────────────────────────────
\section{Experience}

\textbf{Release Engineer} \hfill \textitPubMatic} \hfill \textcolor{mutedtext}{Oct 2025 -- Present}
\begin{itemize}
  \item Managed large-scale weekly production deployments in globally distributed data centers using Jenkins and Ansible.
  \item Designed and optimized CI/ CD pipelines for Private Cloud and AWS, improving deployment reliability and consistency
  \item Built observability stack (Prometheus + Grafana + Loki) reducing MTTR from 45min to 8min
  \item Deployed and managed containerized applications using Docker, Kubernetes, and Helm.
\end{itemize}

\vspace{4pt}

\textbfSystems Engineer} \hfill \textit{PubMatic} \hfill \textcolor{mutedtext}{April 2023 -- Oct 2025}
\begin{itemize}
  \item Automated server provisioning and configuration management using Ansible, improving environment consistency and reducing provisioning time
  \item Developed Bash scripts to automate routine administrative and deployment tasks
  \item Investigated incidents and performed root cause analysis to prevent recurring issues
\end{itemize}

\vspace{4pt}

\textbf{Systems Administrator} \hfill \textit{TCS} \hfill \textcolor{mutedtext}{Jan 2020 -- Nov 2022}
\begin{itemize}
  \item Administered Linux (Ubuntu, CentOS) and Windows servers, handling installation, patching, performance tuning, and security hardening
  \item Automated administrative and monitoring tasks using Bash scripting and Nagios to improve operational efficiency.
  \item Configured networking components including firewalls, load balancers, VPNs, and implemented infrastructure SOP documentation.
\end{itemize}

%── Projects ──────────────────────────────────────────────────
\section{Projects}

\textbf{\href{https://github.com/4-sai/homelab-k8s}{homelab-k8s}} —  Kubernetes cluster with Ansible. Runs full observability stack and self-hosted services.\\[3pt]

%── Certifications ────────────────────────────────────────────
\section{Certifications}

\begin{multicols}{2}
\begin{itemize}
  \item CEH — EC-Council (2019)
\end{itemize}
\end{multicols}

%── Education ─────────────────────────────────────────────────
\section{Education}

\textbf{B.Sc in Computer Science} \hfill \textit{Osmania University} \hfill \textcolor{mutedtext}{2019}

%── CTF / Recognition ─────────────────────────────────────────
\section{CTF \& Recognition}

\begin{itemize}
  \item TryHackMe — Top 5\% globally
\end{itemize}

\end{document}
